\section{n3. New exact separable coefficient-field descent gives all continuous envelopes on finite smooth quotients over any characteristic-zero field.}
\textit{Supports: G331, n1, n2.}

NEW algebraic-field descent, proved independently. Let k be any field, let (R,m) be a positive-dimensional normal domain essentially of finite type over k with residue field EXACTLY k, and assume geometric integrality and geometric normality. Fix primary a. Let k\textasciicircum{}s be a separable closure and put R\_s=R tensor\_k k\textasciicircum{}s. It is the local ring at the corresponding rational point: for every finite Galois E/k, R\_E=R tensor\_k E is finite free and has a unique maximal ideal, since R\_E/mR\_E=E and every maximal ideal in a finite integral algebra lies over m. The directed union is local, and is the corresponding finite-type localized k\textasciicircum{}s algebra, hence Noetherian excellent normal. Geometric hypotheses are explicit; no assertion about arbitrary excellent germs or completions is made. All semivaluation supports are retained.

Write X\_E=V\_(aR\_E)(R\_E), X\_s=V\_(aR\_s)(R\_s), and Gamma=Gal(k\textasciicircum{}s/k). G331 supplies surjective restriction at every finite stage. There is an actual homeomorphism X\_s=lim\_E X\_E: a compatible tuple defines a semivaluation on the union, satisfies the axioms and a-normalization there, and conversely every semivaluation restricts. Coordinates belong to a finite stage, so the topology is exactly the inverse-limit topology. Compactness and finite-stage surjectivity prove rho:X\_s-\textgreater{}X=V\_a(R) is continuous surjective, also at nonminimal supports. One can see surjectivity by the finite-intersection property for the nonempty compact finite-stage fibres.

Every rho-fibre is a FULL Gamma orbit. Indeed R\_E\textasciicircum{}Gal(E/k)=R: expand in a k-basis of E and impose invariance, or use the trace-dual coefficient argument below. Its generic extension is Galois of degree [E:k] by geometric integrality and n2's fixed-field proof. Thus G331's full-support normal-target coefficient trace identifies each finite-stage fibre with a full Gal(E/k) orbit. Given w,w' in the same infinite fibre, let C\_E be the set of gamma in Gamma for which their E-restrictions agree after gamma. Each C\_E is nonempty clopen, because finite-field automorphisms extend to k\textasciicircum{}s and the fibre is an orbit. Finite intersections contain C\_F for a common finite Galois F. The profinite compactness of Gamma gives a common gamma. This proves full transitivity, not merely transitivity on original-field valuations. Gamma is the inverse limit of the finite Galois groups; its compactness and restriction surjectivity are ordinary separable field theory, and the finite-intersection argument is supplied here.

Here is COMPLETE ideal descent, avoiding an infinite average. For finite Galois E/k choose a k-basis e\_j and trace-dual c\_i with Tr(c\_i e\_j)=delta\_ij. Nondegeneracy can be proved without an extra theorem: n2's evaluation-matrix fixed-field proof gives a basis whose matrix M=(sigma(e\_j)) is invertible. The trace Gram matrix is M\textasciicircum{}t M, with determinant (det M)\textasciicircum{}2 nonzero; its entries lie in k. Thus a trace-dual basis exists, also when the characteristic divides the group order. If an ideal I\_E in R\_E is Galois-stable and s=sum\_j r\_j e\_j belongs to it, then sum\_sigma sigma(c\_i s)=r\_i belongs to I\_E intersect R. Hence I\_E=(I\_E intersect R)R\_E.

If an ideal I\_s in R\_s is Gamma-stable, choose a finite Galois E containing its finite set of generators and let I\_E be the generated ideal. R\_s is faithfully flat over R\_E (field extension), so I\_s intersect R\_E=I\_E. Every automorphism of E extends to Gamma, and stability of I\_s forces stability of I\_E. Finite descent therefore gives I\_s=(I\_s intersect R)R\_s. This proves the assertion on the entire Noetherian algebra, not just on a generic field. Primary and unit ideals retain their type: proper primary ideals contain a power of m\_s=mR\_s, and faithful contraction contains a power of m and remains proper.

For any finite continuous signed f on X, put F=f composed with rho. The source canonical B\_m=A\_m\textasciicircum{}(R\_s)(F) is Gamma-stable by its all-support definition. Restriction surjectivity proves B\_m intersect R=A\_m\textasciicircum{}R(f). Ideal descent now gives the exact EVERY-degree identity

 B\_m=A\_m\textasciicircum{}R(f)R\_s,  h\_m\textasciicircum{}(R\_s)=rho\textasciicircum{}*h\_m\textasciicircum{}R,
 direct-sum B\_m t\textasciicircum{}m=(direct-sum A\_m\textasciicircum{}R t\textasciicircum{}m) tensor\_k k\textasciicircum{}s.

Taking the raw supremum identity of n1 proves rho\textasciicircum{}*p\_R(f)=p\_(R\_s)(F), WITHOUT first assuming either envelope continuous. Continuity of either raw envelope is equivalent to continuity of the other: pullback preserves continuity, and continuous fibre-constant functions descend through the compact quotient rho. The uniform criterion of n1 then proves positivity in C whenever continuous. Sup norms and all uniform canonical error bounds are preserved exactly, without a field-degree loss. Finite generation of the two canonical algebras is also equivalent: if the source is generated in degrees \textless{}=D, the multiplication maps from base degrees \textless{}=D become surjective after faithfully flat extension, hence were already surjective; finitely many generators of those Noetherian ideals suffice. The reverse is base extension. This is an equivalence, not general finite generation.

For completeness the positive continuous cone itself descends: rho\textasciicircum{}*C\_R is precisely the Gamma-invariant subclass of C\_(R\_s). Pullback is an isometry and preserves C,D by primary ideal extension and actual limits. If H in source C is invariant, approximate it uniformly within epsilon by -w(I)/r. Descend its generators and primary-power containment to finite Galois E. Replace I\_E by I'\_E=sum\_sigma sigma(I\_E), whose profile is the maximum of its translates and is within epsilon of H on every source point, by Gamma invariance. I'\_E is stable, unit or proper primary. The trace-dual proof descends it to J in R, giving the approximant -v(J)/r downstairs. Compact quotient descent and these uniform approximants put H in C\_R. No assertion that arbitrary invariant D nets automatically descend through an infinite extension is needed or made.

Actual NEW CE application: in characteristic zero k\textasciicircum{}s is an algebraic closure. Let a finite group G act faithfully by k-linear transformations on k[x\_1,...,x\_n], preserving the origin, and take R=(k[x]\textasciicircum{}G) at its origin. The same result holds for any such geometrically normal rational germ whose algebraic-closure ring has the finite smooth-quotient presentation of n2. In the displayed linear case, invariants commute with field extension: the fixed-subspace equations on each finite-dimensional polynomial degree piece commute with tensoring, and localization commutes by invariant orbit denominator products. The invariant ring is finite type and k[x] finite over it: each variable satisfies its monic orbit polynomial, giving finiteness over the finite coefficient subalgebra; the invariant ring is its Noetherian submodule. All nonleading orbit-polynomial coefficients vanish at the origin, so modulo the invariant origin ideal every variable is nilpotent. Thus the origin is the unique centre above it, and the localized smooth cover is actually finite local. Normality follows from the integral invariant-fraction argument in n2, and the same argument after any field extension gives geometric normality and integrality. Thus R\_s is the exact algebraically closed characteristic-zero finite smooth quotient of n2. That note applies BFJ ONLY there. Our independently proved full-fibre and ideal descent now gives a greatest continuous geometric C/D envelope for EVERY finite continuous signed obstacle on R over ANY characteristic-zero field k.

For each signed REAL Cartier obstacle on an ACTUAL primary model, G331's section formula is a finite signed primary combination. Its scalar pullback is therefore a model obstacle in n2. The n2 quotient O(1/m) full-degree canonical estimate transfers EXACTLY to R through the every-degree identity above. The constant may depend on the geometric source model and obstacle; there is no field-extension-degree factor or finite-generation hypothesis. Exact primary-gauge changes leave all canonical ideals unchanged and divide errors by the positive gauge denominator.

In particular k=Q and R=Q[a,b,c,z\_1,...,z\_r]\_(a,b,c,z)/(ac-b\textasciicircum{}2), r\textgreater{}=1, are actual normal nonisolated germs covered by the theorem. Normality, dimension r+2 and singular locus of dimension r follow from n2's explicit invariant presentation and derivative/embedding-dimension computation, which remain valid over characteristic-zero k. The same nonlinear signed primary obstacle and two cancellation-sensitive valuations in n2 exist over Q. Its continuous geometric envelope now follows over Q itself, rather than by applying a published algebraically closed theorem directly to Q. This supplies a complete new scoped theorem beyond whole410's stated base-field domain. It does not treat arbitrary normal singularities without the geometric smooth-quotient presentation, mixed characteristic, or provide a positive-characteristic CE source theorem. All normalization/conductor rules remain conditional on EVERY actual fibre; products and energy remain separate.

Artifacts: n3.txt, distinct full TeX exposition n3.tex, P12.3-algebraic-descent-envelopes.tex and the entire companion/reading ledger. Rerun: python assemble\_delivery.py; python finite\_checks.py; python check\_delivery.py. The universal arguments are proofs; the trace-dual finite example checks only its stated arithmetic domain.

\subsection*{Expanded field-descent equalities}
The infinite group is used by finite clopen compatibility conditions:
\[
 X_{k^s}=\varprojlim_E X_E,\qquad
 \bigcap_E C_E\ne\varnothing,\qquad
 \rho^{-1}(v)=\Gamma\cdot w.
\]
Trace-dual extraction on a stable ideal is an actual ring identity:
\[
 s=\sum_jr_je_j\in I_E\ \Longrightarrow\
 r_i=\sum_{\sigma\in\operatorname{Gal}(E/k)}\sigma(c_is)
 \in I_E\cap R.
\]
It uses separability, not division by the group order. In particular it
also applies when the field characteristic divides that order. It yields
\[
 A_m^{R_{k^s}}(\rho^*f)=A_m^R(f)R_{k^s},\qquad
 h_m^{R_{k^s}}=\rho^*h_m^R,\qquad
 p_{R_{k^s}}(\rho^*f)=\rho^*p_R(f).
\]
Unlike a coefficient norm approximation, this carries the canonical
degree $m$ itself unchanged. The characteristic-zero application is a
descent theorem proved here, after using the published smooth theory
only over the algebraic closure. It is not an enlarged citation of BFJ.

